
\begin{obereflection}
\textbf{Latihan 3.1}
\begin{enumerate}
    \item Hitunglah nilai dari $-24 \pmod{11}$.
    \item Tentukan invers modulo dari $12 \pmod{5}$.
    \item Diketahui sebuah pesan memiliki huruf A, B, C, D dengan probabilitas
          masing-masing $p(A)=0{,}25$, $p(B)=0{,}25$, $p(C)=0{,}25$, $p(D)=0{,}25$.
          Hitunglah entropi pesan tersebut.
    \item Tentukan $\text{PBB}(48, 18)$ dengan Algoritma Euclidean, lalu nyatakan
          hasilnya sebagai kombinasi lanjar $m \cdot 48 + n \cdot 18$.
    \item Hitunglah $\varphi(20)$, $\varphi(7)$, dan $\varphi(8)$. Jelaskan mengapa
          hasil $\varphi(20)$ sama dengan $\varphi(4)\cdot\varphi(5)$.
    \item Periksa Teorema Euler untuk $a = 3$ dan $n = 10$: hitung $\varphi(10)$ dan
          tunjukkan bahwa $3^{\varphi(10)} \equiv 1 \pmod{10}$.
    \item Tunjukkan bahwa $g = 3$ adalah akar primitif modulo 7, dan bahwa $g = 2$
          bukan. Tuliskan seluruh pangkatnya sampai kembali ke 1.
    \item Diketahui $7^x \equiv 15 \pmod{41}$. Carilah $x$. Jelaskan mengapa
          menghitung $7^x$ dari $x$ yang diketahui jauh lebih mudah daripada
          menentukan $x$ dari hasilnya.
    \item Sebuah pesan 200 huruf memiliki entropi 4,2 bit/karakter. Berapa persen
          entropi maksimum alfabet 26 huruf yang dicapainya, dan apa artinya bagi
          kekuatan cipherteks tersebut terhadap analisis frekuensi?
\end{enumerate}
\end{obereflection}
